\scanpage{92}
\mspage{98}

des objets localement constants sur $Z$ s'exprime en termes du $(i+1)$-type d'homotopie, mais pas en termes de son $i$-type d'homotopie \unsure{seulement}.

Il semble bien que ce soit l'axiome \uline{Der 3} (a') qui assure l'existence d'une bonne notion \inserted{(``\unsure{bilatér.|bilatère}'')} de $\mathbb D$-fibration, qui soit stable par changement de base, et qui se réduise soit à ``$\mathbb D$-fibration à gauche universelle'', soit à ``$\mathbb D$-fibration à droite universelle'' --- donc qui soit en plus \inserted{une condition} \uline{vérifiable} pratiquement.

Dans le cas où $f$ est propre ou lisse, le critère est particulièrement simple, puisqu'il se borne à la condition que \deleted{\unsure{pour}} les $f_*f^*\deleted{\unsure{A}}a_Y$, resp\textsuperscript{t.} les $f_!f^*a_Y$ \inserted{(pour $a\in A$)}, soient des objets loc. constants sur $Y$ (donc \uline{constants} si $Y$ est \uuline{Hot}-asphérique, en vertu de \unsure{Der,|Der} 3 a').

Je vais également introduire une variante \deleted{\unsure{de}} forte de \uline{Der} 3 b) :

% slanted
\margin{\uline{Der 3} $b'_0$ : le foncteur est fidèle \unsure{(??)|(?)} $b'_1$ : le foncteur est pl. fidèle $b_2$ : le foncteur \unsure{est} une équiv.}

\uline{Der 3} b') Soit $Z'\unsure{\relbar|\to}Z$ une $\mathrm{Hot}$-équivalence dans $\mathrm{Diag}$. Alors $f^*:\mathbb D(Z')\unsure{\relbar|\to}\mathbb D(\unsure{Z|2})$ induit une équivalence entre les catégories des \deleted{\unsure{objets}}coeff. \uline{loc. constants} sur $Z$ et sur $Z'$.

% archive sheet number: 91
